% WORSTCASE REPORT::BRIDGE_AUDIT1::2026-11-10
\documentclass{article}\usepackage{amsmath,amssymb,amsthm}
\begin{document}
\title{Worst-case derivation cost: audit of the bridge $\mathrm{SAT}\in P\Rightarrow \overline W_R\le\mathrm{poly}$}
\maketitle
Status of everything below: P vs NP OPEN; no lower bound on $\overline W_R$ for a strong $R$ is claimed.

\section*{Setup}
$R$ is a sound proof system with polynomial-time verifier $V_R$; $\mathrm{cost}_R(\pi)=|\pi|+$ verifier time.
$\overline W_R(N)=\max_{F\in\mathrm{UNSAT},|F|\le N}\min_\pi \mathrm{cost}_R(\pi)$. Free choice of $\pi$ is built in.

\section*{Results (all KNOWN or elementary, restated to fix quantifiers)}
\paragraph{T1 (KNOWN, Cook--Reckhow).} $\overline W_R\le \mathrm{poly}$ for some sound $R$ iff NP$=$coNP. Hence ``$\overline W_R$ superpolynomial for \emph{all} sound $R$'' is equivalent to NP$\ne$coNP, strictly stronger than P$\ne$NP.
\paragraph{T2 (PROVED/EXACT, trivial).} If SAT$\in$P via $M$, then $R_M=\{(F,\varepsilon): M(F)\text{ rejects}\}$ is sound and polynomially bounded. So an \emph{existential} bridge is trivial and useless: it does not constrain a fixed natural $R$.
\paragraph{T3 (reformulation, PROVED/EXACT).} For fixed $R$, define $B(R)$: $P=NP\Rightarrow \overline W_R\le\mathrm{poly}$. Then $B(R)\iff (P\neq NP)\lor(\overline W_R\le \mathrm{poly})$. Any proof of $B(R)$ not proving $P\ne NP$ must build, from a hypothetical algorithm $M$, $R$-proofs while $R$ (hence $V_R$) is fixed \emph{before} $M$. $V_R$ cannot trust $M$, so the proof must certify the correctness of $M$ on $F$ inside $R$: a reflection-type obligation.
\paragraph{T4 (KNOWN, Kraj\'i\v cek--Pudl\'ak 1989; Messner 1999).} E$=$NE implies a p-optimal system exists, and P$=$NP implies E$=$NE by padding. So \emph{in the world P$=$NP} some $R^\*$ satisfies $B(R^\*)$, but $R^\*$ is not explicitly defined unconditionally. Whether $\mathrm{TAUT}$ has a p-optimal/optimal system unconditionally is OPEN.
\paragraph{Consequence (CANDIDATE reading).} The ``weakest sufficient bridge'' for a \emph{natural} $R$ (Frege, EF, resolution$+$XOR$+$CP$+$PC combos) is, in its uniform form, at least as hard as the optimal-proof-system problem restricted to $R$: it holds for EF if the SAT algorithm's correctness is provable in $PV$ (Cook), but provability is not implied by $SAT\in P$. So the bridge needs either (a) a provability hypothesis (CONDITIONAL), or (b) an $R$ with built-in reflection, whose lower bounds are then as hard as T1.

\section*{Where lower bounds are actually known (literature check, 2026-11-10)}
Resolution, CP, PC, bounded-depth Frege: exponential for random $k$-CNF / PHP / Tseitin (KNOWN). Res$(\oplus)$/ResLin: Tseitin lifted with IP gadget needs $\exp(\tilde\Omega(N^\epsilon))$ for depth $O(N^{2-\epsilon})$ (Bhattacharya--Chattopadhyay, arXiv 2507.23008); random-CNF ResLin bound OPEN. Res$(k)$ random CNF: Sofronova--Sokolov. For a ``combined system'' the correct next target is thus ResLin on random lifted CNF; nothing is known for Frege.

\section*{Red-team notes}
Hard family candidates must survive: XOR (Tseitin falls to GF(2)), extension variables, shared DAG (tree-size vs DAG-size), adaptive switching between rules. Tseitin and PHP-type families fail against PC/CP; random 3-CNF survives all weak systems but no lower bound is known for Frege.

\section*{Next targets}
(i) Formalize multisigned DAG hypergraph states and cost; (ii) decide whether a restricted reflection rule $R^{\mathrm{refl}}_k$ (reflection for algorithms of size $\le k$) gives $B$ for the sub-bridge and what T1-type cost it brings; (iii) ResLin random-CNF lower bound.
% END_OF_VERIFIED_REPORT::WORSTCASE::BRIDGE_AUDIT1
\end{document}
